%% =====================================================================
%%  Corrigés — Unit 12 « I can swim! »
%% =====================================================================
\corrigesUnite{12}{I can swim!}

\begin{corrige}{1}
\begin{reponses}[4]
\item \rep{swim}
\item \rep{ride a bike}
\item \rep{play the piano}
\item \rep{drive a car}
\item \rep{sing}
\item \rep{play tennis}
\item \rep{draw}
\item \rep{run}
\end{reponses}
\retiens{\en{play the piano} (instrument : \en{the}) mais \en{play tennis} (sport : pas d'article).}
\end{corrige}

\begin{corrige}{2}
\begin{reponses}[4]
\item \rep{can}
\item \rep{can't}
\item \rep{can't}
\item \rep{can}
\item \rep{can't}
\item \rep{can}
\item \rep{can't}
\item \rep{can}
\end{reponses}
\end{corrige}

\begin{corrige}{3}
\begin{reponses}[2]
\item \rep{Yes, she can.}
\item \rep{No, he can't.}
\item \rep{No, she can't.}
\item \rep{Yes, he can.}
\item \rep{Yes, they can.}
\end{reponses}
\pourquoi{la réponse courte reprend \en{can} avec le bon pronom : Léa → \en{she}, Jack → \en{he},
Léa et Emma → \en{they}.}
\end{corrige}

\begin{corrige}{4}
\begin{reponses}[3]
\item \rep{very well}
\item \rep{quite well}
\item \rep{not very well}
\item \rep{at all}
\item \rep{very fast}
\end{reponses}
\pourquoi{\en{at all} renforce une négation : \en{I can't ski at all} = « je ne sais pas skier du tout ».}
\end{corrige}

\begin{corrige}{5}
\begin{reponses}[1]
\item \rep{She can swim.}
\item \rep{I can drive.}
\item \rep{Can you ski?}
\item \rep{He sings very well.}
\item \rep{He can't dance.}
\end{reponses}
\pourquoi{\en{can} n'a ni \en{-s}, ni \en{to}, ni \en{do}. Phrase 4 : un adverbe (\en{well}) décrit le
verbe \en{sings}.}
\end{corrige}

\begin{corrige}{6}
\begin{reponses}[3][french]
\item \rep{V}
\item \rep{F}
\item \rep{V}
\item \rep{F}
\item \rep{F}
\item \rep{V}
\end{reponses}
\pourquoi{Kim joue du piano (Ben est à la batterie) ; le concert a lieu vendredi ; on n'a pas besoin
de lire la musique (\en{You don't need to read music}).}
\end{corrige}

\begin{corrige}{7}
\begin{reponses}[3]
\item \rep{can}
\item \rep{can't}
\item \rep{can't}
\item \rep{can}
\item \rep{can't}
\item \rep{can}
\end{reponses}
\pourquoi{dans les phrases avec \en{can't}, le verbe auxiliaire est accentué et le \phon{æ} bien ouvert ;
avec \en{can}, il est faible et l'accent tombe sur le verbe principal. Texte de la piste
page~\pageref{trans-12.10}.}
\end{corrige}

\begin{corrige}{8}
\begin{reponses}[5]
\item \rep{yes}
\item \rep{no}
\item \rep{yes}
\item \rep{no}
\item \rep{yes}
\end{reponses}
\pourquoi{Maya parle espagnol : sa grand-mère vient de Porto Rico (\en{Puerto Rico}).}
\end{corrige}

\begin{corrige}{9}
\begin{reponses}[4]
\item \rep{Can}
\item \rep{can}
\item \rep{well}
\item \rep{swim}
\item \rep{can't}
\item \rep{all}
\item \rep{run}
\item \rep{fast}
\end{reponses}
\end{corrige}

\begin{corrige}{10}
\begin{reponses}[1]
\item \rep{I can swim, but I can't ski.}
\item \rep{Can you play the guitar? — Yes, I can.}
\item \rep{She sings very well.}
\item \rep{He can't drive.}
\item \rep{Can you come on Friday?}
\item \rep{I play the piano, but not very well.}
\item \rep{They can't dance at all.}
\item \rep{She's a good swimmer.}
\end{reponses}
\piege{phrase 8 : \en{good} (adjectif) devant le nom \en{swimmer}, et non \en{well}.}
\end{corrige}

\begin{corrige}{11}
Réponse modèle :
\begin{anglais}
\begin{tcolorbox}[otBase,colback=otGreenTint,colframe=otGreenTint,fontupper=\fontsize{10.4}{14}\selectfont]
I can swim very well and I can ride a bike. I can cook quite well. I can play the guitar, but not
very well. I can't ski at all, and I can't sing! My sister can sing very well.
\end{tcolorbox}
\end{anglais}
\end{corrige}

\begin{corrige}{12}
Réponses modèles :
\begin{reponses}[1]
\item \rep{Yes, I can. I swim quite well.}
\item \rep{No, I can't.}
\item \rep{I can play tennis and I can run.}
\item \rep{Yes, I can. I can play the piano.} \textit{ou} \rep{No, I can't.}
\item \rep{Yes, I can!}
\end{reponses}
\end{corrige}

\begin{corrigetest}
\begin{reponses}[5][french]
\item \rep{b}
\item \rep{c}
\item \rep{a}
\item \rep{b}
\item \rep{a}
\item \rep{c}
\item \rep{b}
\item \rep{a}
\item \rep{c}
\item \rep{b}
\end{reponses}
\pourquoi{question 4 : \en{well} est l'adverbe de \en{good}. Question 6 : pas d'article devant un sport.}
\end{corrigetest}
